%% Poids des bits : décomposition de (1101)_2 = 8 + 4 + 0 + 1 = 13.
%% Le poids d'un bit dépend de sa position ; il double de la droite vers la gauche.
\documentclass[tikz,border=4pt]{standalone}
\usepackage{liaisons}
\begin{document}
\begin{tikzpicture}[x=1cm,y=1cm,
  case/.style={draw=solideB, line width=1.1pt, fill=white, minimum size=0.7cm,
               inner sep=0pt, font=\small},
  ligne/.style={font=\scriptsize, anchor=east},
  equerre/.style={solideB, line width=0.7pt, -{Stealth[length=4pt]}, rounded corners=2pt},
  fl/.style={-{Stealth[length=4.5pt]}, line width=0.7pt, solideA},
  note/.style={font=\scriptsize, text=solideA, text width=1.5cm}]

%% ================== les trois lignes ==================================
\foreach \i/\p/\v/\c in {0/3/{2^3}/1, 1/2/{2^2}/1, 2/1/{2^1}/0, 3/0/{2^0}/1} {
  \node[font=\scriptsize] at (\i*0.8,1.3) {\p};
  \node[font=\scriptsize] at (\i*0.8,0.75) {$\v$};
  \node[case] (c\i) at (\i*0.8,0) {\c}; }
\node[ligne] at (-0.55,1.3)  {Position};
\node[ligne] at (-0.55,0.75) {Valeur};
\node[ligne] at (-0.55,0)    {Chiffres};

%% ================== poids fort / poids faible =========================
\node[note, align=right, anchor=east] at (-0.85,-0.95) {bit de poids fort (msb)};
\draw[fl] (-0.8,-0.85) -- (-0.32,-0.3);
\node[note, align=left, anchor=west] at (3.3,0) {bit de poids faible (lsb)};
\draw[fl] (3.25,0) -- (2.8,0);

%% ================== développement en somme de poids ===================
\foreach \i/\y/\t in {3/-0.6/{$1 \times 2^0 = 1$}, 2/-1.05/{$0 \times 2^1 = 0$},
                      1/-1.5/{$1 \times 2^2 = 4$}, 0/-1.95/{$1 \times 2^3 = 8$}} {
  \draw[equerre] (\i*0.8,-0.35) -- (\i*0.8,\y) -- (4.05,\y);
  \node[font=\scriptsize, anchor=west, inner sep=2pt] at (4.15,\y) {\t}; }
\draw[black!35, line width=0.6pt, rounded corners=3pt] (4.1,-2.2) rectangle (5.85,-0.38);
\node[draw=solideE, line width=1pt, rounded corners=3pt, fill=solideE!12,
      font=\small\bfseries, inner sep=3pt, anchor=north west] at (4.1,-2.4) {Total = 13};
\end{tikzpicture}
\end{document}
